% discussion.tex -- Sections V and VI. Sources as in results.tex. Facts about Chong et al.:
% the published text describes SO+JR labels and states that M_add > 0 loads the medial
% compartment (paper/author.tex lines 537, 554); its reductions were not checked with SO+JR
% (ISRR2). Kaneda 0.28 BW = our recomputation on SC (gc/pilot_AB_sojr/README.md); "5 of 94
% segments cover a full stance" (qc_final.md fix:label-qc open issues).
% "Cannot confirm or refute" rests on D1 and D2, which hold for every checkpoint trained on the stored
% labels and for every optimizer run that forms the proxy with +M_x at the right knee, NOT on the in
% vivo r of the two audited checkpoints: which checkpoint produced the figures of [4] is unconfirmed
% (facts ISRR4; TODO before submission: confirm with the first author). SC residual over-read:
% audit_gc/qc_keep/summary.csv; 25.4 %: labels_qc/label_masks/by_lab.csv (see results.tex).
% Moore share: 92.6 % of QC-passing training frames (labels_qc/label_masks/dry_build_fpopt/
% kept_by_lab.csv, 327,287 / 353,581) and 92.9 % of held-out stance frames (fmedial_eval/
% audit_test_qc/summary.csv, 31,515 / 33,915). Optimizer input convention for the pelvis angles is
% unverified (sep23/README.md; evals_final.md fix:sep23). The reference motion's COP is pinned at the optimizer's
% contact point, which alone fails the fixed-COP rule (sep23/README.md), so verifying synthesized gaits needs a
% moving COP as well as converted pelvis angles. Test split: one person per lab (fmedial_eval/audit_test_qc).
\section{Discussion}\label{sec:discussion}
\textit{What the audit means for the optimizer.} \Chong{}~\citeChong{} describe SO+JR labels
and a medial-positive moment, but the stored labels lack the SO muscle forces (D1) and the
inference code uses the inverted sign (D2). D1 affects every checkpoint trained on the stored
labels and D2 every optimizer run with this proxy, whichever checkpoint it queried; held-out
accuracy hid both, because training and test labels share them. The reported load reductions are
surrogate predictions not checked with SO+JR, so this audit can neither confirm nor refute them,
and the in vivo scores of the two audited checkpoints need not carry over to others. On the
optimizer's reference motion, corrected and legacy models fall on opposite sides of the limit.
Running SO+JR on the synthesized gaits, with converted pelvis angles and a moving COP, would test in
the model whether they unload the knee; in vivo confirmation needs measurement.

\textit{Neither label source is ground truth.} The stored labels are too low and the SO+JR
labels too high. The quality rules remove the worst of D3 but keep only a quarter of stance time
and little of Carter or Lencioni, and frames that pass them are still over-read in all three knees
(Sec.~\ref{sec:res-d3}); the corrected baseline inherits this. On SC, Kaneda et al.'s SO+JR, with
separate medial and lateral contacts, fitted the implant far better (Sec.~\ref{sec:res-invivo}),
so SO+JR can do better than our labels; their contact model may explain the difference.
Re-solving SO per frame and repairing the COP would keep more data.

\textit{Recommendations.} For learned models whose outputs act on people: (i)~check that the physics ran
as described (D1); (ii)~test sign and axis conventions on a case with a known answer (D2, S2);
(iii)~compare label magnitudes with in vivo data (D3); (iv)~build features with the deployment
code (D4); (v)~report seed spread off the training distribution, and verify optimized motions
with physics before they reach a patient.

\textit{Social impact.} Gait retraining has lowered KAM and pain~\cite{shull2013six}, and a
wearable surrogate could take pain-aware coaching out of the lab once contact detection works in
slow walking and its loads agree with in vivo measurement. The main risk is false reassurance: an under-reading model certifies a harmful gait as
safe, while an over-reading one steers patients toward needless changes. $\Fmed$ remains a proxy, and medial load is not pain, so patient use
also needs clinical validation.

\textit{Limitations.} The in vivo subjects are three knee-replacement patients, not OA knees; most
segments cover only half a stance, peak errors rest on 3--6 windows per knee and half, and DM's
timing is uncertain. Our SO+JR uses a single knee contact. The 40 peaks were hand-picked, so their
causes are not prevalences. The test split has one person per lab; Moore supplies 93\% of the
quality-checked training frames, and one Moore person 93\% of the held-out frames. The optimizer
was not available, so its pelvis-angle convention is unverified, and S1 is not propagated through
it (future work). The real-sensor data have no knee-load labels, and the suit's label error rests
on one trial.

\section{Conclusion}\label{sec:conclusion}
A learned pain-risk surrogate inherits the defects of its label pipeline and input construction,
and held-out accuracy hides those that the test data share. Auditing one such pipeline against
physics, instrumented knees and its own code found four defects. Retrained on quality-checked
labels with the physical sign, the surrogate correlates positively with the implants but less well
than its labels; on the optimizer's reference motion its seeds agree on the verdict but differ by
up to 1.0~BW in load. A load reduction found by optimizing such a
surrogate is a hypothesis to verify with physics.
